\documentclass[14pt,a4paper]{extarticle}

\usepackage[T2A]{fontenc}
\usepackage[utf8]{inputenc}
\usepackage[russian]{babel}
\usepackage{cmap}
\usepackage{geometry}
\usepackage{array}
\usepackage{longtable}
\usepackage{tabularx}
\usepackage{booktabs}
\usepackage{xcolor}
\usepackage{ragged2e}
\usepackage{seqsplit}
\usepackage{listings}
\usepackage{enumitem}
\usepackage{hyperref}
\usepackage{tikz}
\usetikzlibrary{arrows.meta,positioning,shapes.geometric}

\geometry{
    left=30mm,
    right=15mm,
    top=20mm,
    bottom=20mm
}

\hypersetup{
    colorlinks=true,
    linkcolor=black,
    urlcolor=blue
}

\setlength{\parindent}{1.25cm}
\setlength{\parskip}{0.2em}
\renewcommand{\arraystretch}{1.35}
\setlength{\tabcolsep}{4pt}
\sloppy

\newcolumntype{L}[1]{>{\RaggedRight\arraybackslash}p{#1}}

\lstdefinestyle{jsonstyle}{
    basicstyle=\ttfamily\small,
    breaklines=true,
    frame=single,
    columns=fullflexible,
    keepspaces=true,
    showstringspaces=false,
    backgroundcolor=\color{gray!5}
}

\newcommand{\code}[1]{\texttt{#1}}
\newcommand{\pathcode}[1]{{\footnotesize\ttfamily\seqsplit{#1}}}

\begin{document}

\begin{titlepage}
    \begin{center}
        САНКТ-ПЕТЕРБУРГСКИЙ НАЦИОНАЛЬНЫЙ\\
        ИССЛЕДОВАТЕЛЬСКИЙ УНИВЕРСИТЕТ ИТМО

        \vspace{1.5cm}

        Дисциплина: Бэк-энд разработка

        \vspace{1.5cm}

        \textbf{Отчет}

        \vspace{0.5cm}

        Домашнее задание 4

        \vspace{0.5cm}

        \textbf{Технический дизайн микросервисной архитектуры}

        \vfill

        \begin{flushright}
            Выполнил:\\
            Цатинян Артём\\
            Группа БР1.2

            \vspace{0.8cm}

            Проверил:\\
            Добряков Д. И.
        \end{flushright}

        \vfill

        Санкт-Петербург\\
        2026 г.
    \end{center}
\end{titlepage}

\section{Задача}

Тема проекта: приложение для бронирования столиков в ресторанах.

В рамках домашнего задания требовалось спроектировать переход от монолитного backend-приложения к микросервисной архитектуре. Необходимо разделить приложение на независимые сервисы по принципу \code{database-per-service}, описать взаимодействие между сервисами, спроектировать разделение существующей базы данных, подготовить OpenAPI-спецификацию для межсервисного взаимодействия, а также описать запросы, ответы, возможные ошибки и шаги миграции.

\section{Ход работы}

Исходное приложение реализовано как монолитный REST API на Kotlin и Spring Boot. В одной базе данных находятся пользователи, рестораны, кухни, столики, меню, бронирования и отзывы. Такая структура удобна для первой реализации, но в монолите разные части предметной области связаны прямыми SQL-запросами и внешними ключами. При переходе к микросервисам эти связи нужно заменить на четкие контракты между сервисами.

Декомпозиция выполняется не по каждой таблице отдельно, а по самостоятельным бизнес-областям. Это позволяет не переусложнять архитектуру и при этом соблюсти принцип владения данными.

\subsection{Выделенные сервисы}

\begin{longtable}{|L{0.23\textwidth}|L{0.47\textwidth}|L{0.18\textwidth}|}
\hline
\textbf{Сервис} & \textbf{Ответственность} & \textbf{Собственная БД} \\
\hline
API Gateway & Единая внешняя точка входа, маршрутизация публичных запросов, проверка JWT, прокидывание \code{userId} в downstream-сервисы & Нет \\
\hline
Identity Service & Регистрация, вход, хранение пользователей, ролей, паролей, профиля пользователя & \code{identity\_db} \\
\hline
Restaurant Catalog Service & Каталог ресторанов, кухни, адреса, рабочие часы, столики, фотографии, меню, агрегированные рейтинги & \code{catalog\_db} \\
\hline
Booking Service & Создание бронирований, отмена, просмотр бронирований, расчет доступности столиков & \code{booking\_db} \\
\hline
Review Service & Создание и просмотр отзывов, проверка права пользователя оставить отзыв, события изменения рейтинга & \code{review\_db} \\
\hline
\end{longtable}

API Gateway не владеет бизнес-данными. Его задача --- сохранить для клиента прежний внешний REST API и скрыть от клиента внутреннее разбиение системы.

\subsection{Архитектурная схема}

\begin{center}
\begin{tikzpicture}[
    service/.style={rectangle, draw, rounded corners, align=center, text width=3.1cm, minimum height=0.9cm, font=\small},
    db/.style={cylinder, draw, shape border rotate=90, aspect=0.25, align=center, text width=2.3cm, minimum height=0.85cm, font=\small},
    arrow/.style={-{Latex}, thick}
]
    \node[service] (client) at (0, 0) {Client\\Postman / Web};
    \node[service] (gateway) at (0, -1.5) {API Gateway};

    \node[service] (identity) at (4.6, 0.1) {Identity\\Service};
    \node[service] (catalog) at (4.6, -1.3) {Restaurant Catalog\\Service};
    \node[service] (booking) at (4.6, -2.7) {Booking\\Service};
    \node[service] (review) at (4.6, -4.1) {Review\\Service};

    \node[db] (identitydb) at (8.4, 0.1) {\code{identity\_db}};
    \node[db] (catalogdb) at (8.4, -1.3) {\code{catalog\_db}};
    \node[db] (bookingdb) at (8.4, -2.7) {\code{booking\_db}};
    \node[db] (reviewdb) at (8.4, -4.1) {\code{review\_db}};

    \draw[arrow] (client) -- (gateway);
    \draw[arrow] (gateway.east) -- (identity.west);
    \draw[arrow] (gateway.east) -- (catalog.west);
    \draw[arrow] (gateway.east) -- (booking.west);
    \draw[arrow] (gateway.east) -- (review.west);

    \draw[arrow] (identity) -- (identitydb);
    \draw[arrow] (catalog) -- (catalogdb);
    \draw[arrow] (booking) -- (bookingdb);
    \draw[arrow] (review) -- (reviewdb);
\end{tikzpicture}
\end{center}

На схеме показаны внешняя маршрутизация и владение базами данных. Межсервисные вызовы вынесены в отдельную таблицу ниже, чтобы не перегружать диаграмму пересекающимися стрелками. Основной тип синхронного взаимодействия между сервисами --- REST по внутренним эндпоинтам. Для обновления агрегированного рейтинга ресторана используется асинхронное событие \code{ReviewCreated}, потому что рейтинг не должен блокировать создание отзыва распределенной транзакцией.

\subsection{Маршрутизация внешнего API}

Внешние пути можно оставить совместимыми с текущим API, чтобы клиентам не пришлось знать о внутренней декомпозиции.

\begin{longtable}{|L{0.60\textwidth}|L{0.30\textwidth}|}
\hline
\textbf{Внешний endpoint} & \textbf{Целевой сервис} \\
\hline
\pathcode{POST /api/v1/auth/register} & Identity Service \\
\hline
\pathcode{POST /api/v1/auth/login} & Identity Service \\
\hline
\pathcode{GET /api/v1/users/me}, \pathcode{PATCH /api/v1/users/me} & Identity Service \\
\hline
\pathcode{GET /api/v1/cuisines} & Restaurant Catalog Service \\
\hline
\pathcode{GET /api/v1/restaurants} & Restaurant Catalog Service \\
\hline
\pathcode{GET /api/v1/restaurants/\{restaurantId\}} & Restaurant Catalog Service \\
\hline
\pathcode{GET /api/v1/restaurants/\{restaurantId\}/menu} & Restaurant Catalog Service \\
\hline
\pathcode{GET /api/v1/restaurants/\{restaurantId\}/availability} & Booking Service \\
\hline
\pathcode{POST /api/v1/bookings} & Booking Service \\
\hline
\pathcode{GET /api/v1/bookings/me} & Booking Service \\
\hline
\pathcode{GET /api/v1/bookings/\{bookingId\}} & Booking Service \\
\hline
\pathcode{PATCH /api/v1/bookings/\{bookingId\}/cancel} & Booking Service \\
\hline
\pathcode{GET /api/v1/restaurants/\{restaurantId\}/reviews} & Review Service \\
\hline
\pathcode{POST /api/v1/restaurants/\{restaurantId\}/reviews} & Review Service \\
\hline
\end{longtable}

Endpoint доступности столиков относится к Booking Service, потому что именно этот сервис владеет бронированиями и может корректно определить пересечения по времени. При этом Booking Service получает из Restaurant Catalog Service список активных столиков, вместимость и рабочие часы ресторана.

\subsection{Разделение базы данных}

В микросервисной архитектуре сервисы не обращаются напрямую к чужим таблицам. Внешние ключи между разными бизнес-областями заменяются обычными идентификаторами и проверками через внутренние API.

\subsubsection{\code{identity\_db}}

\begin{longtable}{|L{0.28\textwidth}|L{0.62\textwidth}|}
\hline
\textbf{Таблица} & \textbf{Назначение} \\
\hline
\code{users} & Пользователи, email, хэш пароля, имя, фамилия, телефон, роль, активность, даты создания и обновления \\
\hline
\end{longtable}

Identity Service является единственным владельцем учетных данных. Остальные сервисы хранят только \code{user\_id} как внешний идентификатор и при необходимости получают краткую информацию о пользователе через internal API.

\subsubsection{\code{catalog\_db}}

\begin{longtable}{|L{0.32\textwidth}|L{0.58\textwidth}|}
\hline
\textbf{Таблица} & \textbf{Назначение} \\
\hline
\code{restaurants} & Основная информация о ресторанах, адрес, контакты, ценовой сегмент, политика бронирования \\
\hline
\code{cuisines} & Справочник кухонь \\
\hline
\code{restaurant\_cuisines} & Связь ресторанов и кухонь \\
\hline
\code{restaurant\_working\_hours} & Рабочие часы ресторанов \\
\hline
\code{restaurant\_tables} & Столики ресторанов и вместимость \\
\hline
\code{restaurant\_photos} & Фотографии ресторанов \\
\hline
\code{menu\_categories} & Категории меню \\
\hline
\code{menu\_items} & Позиции меню \\
\hline
\code{restaurant\_rating\_stats} & Денормализованный read-model: средний рейтинг и количество отзывов \\
\hline
\end{longtable}

Таблица \code{restaurant\_rating\_stats} нужна для сортировки и фильтрации ресторанов по рейтингу без синхронных запросов в Review Service на каждый поиск.

\subsubsection{\code{booking\_db}}

\begin{longtable}{|L{0.28\textwidth}|L{0.62\textwidth}|}
\hline
\textbf{Таблица} & \textbf{Назначение} \\
\hline
\code{bookings} & Бронирования пользователя: \code{user\_id}, \code{restaurant\_id}, \code{table\_id}, статус, период, количество гостей, пожелания \\
\hline
\end{longtable}

В \code{bookings} больше нет внешних ключей на \code{users}, \code{restaurants} и \code{restaurant\_tables}, потому что эти таблицы принадлежат другим сервисам. Для удобства отображения можно добавить snapshot-поля: \code{restaurant\_name\_snapshot}, \code{table\_number\_snapshot}, \code{table\_seats\_snapshot}. Они фиксируют состояние на момент создания бронирования и позволяют отдавать список бронирований без лишних синхронных вызовов.

\subsubsection{\code{review\_db}}

\begin{longtable}{|L{0.28\textwidth}|L{0.62\textwidth}|}
\hline
\textbf{Таблица} & \textbf{Назначение} \\
\hline
\code{reviews} & Отзывы: \code{booking\_id}, \code{restaurant\_id}, \code{user\_id}, оценка, комментарий, snapshot имени автора \\
\hline
\code{review\_outbox} & Очередь исходящих событий для надежной доставки \code{ReviewCreated} \\
\hline
\end{longtable}

В \code{reviews.booking\_id} сохраняется уникальное ограничение внутри \code{review\_db}, чтобы по одному бронированию нельзя было оставить два отзыва. Проверка того, что бронирование действительно принадлежит пользователю и завершено, выполняется через Booking Service.

\subsection{Замена связей монолита}

\begin{longtable}{|L{0.38\textwidth}|L{0.52\textwidth}|}
\hline
\textbf{Связь в монолите} & \textbf{Решение после декомпозиции} \\
\hline
\code{bookings.user\_id -> users.id} & \code{booking\_db.bookings.user\_id} хранится как scalar id, пользователь проверяется через JWT и при необходимости через Identity Service \\
\hline
\code{bookings.restaurant\_id -> restaurants.id} & Booking Service вызывает Restaurant Catalog Service перед созданием бронирования \\
\hline
\code{bookings.table\_id -> restaurant\_tables.id} & Booking Service получает table context из Restaurant Catalog Service и проверяет вместимость \\
\hline
\code{reviews.booking\_id -> bookings.id} & Review Service вызывает Booking Service, чтобы проверить завершенное бронирование \\
\hline
\code{reviews.user\_id -> users.id} & Review Service получает имя автора из JWT или Identity Service и хранит snapshot \\
\hline
\code{restaurants} joins \code{reviews} для рейтинга & Catalog Service читает локальную таблицу \code{restaurant\_rating\_stats}, обновляемую событием из Review Service \\
\hline
\code{restaurant\_tables} joins \code{bookings} для availability & Booking Service считает доступность по своей БД бронирований и данным о столиках из Catalog Service \\
\hline
\end{longtable}

\subsection{Межсервисное взаимодействие}

OpenAPI-спецификация внутренних REST-контрактов вынесена в файл \code{homeworks/hw4/interservice-openapi.yaml}.

\begin{longtable}{|L{0.46\textwidth}|L{0.21\textwidth}|L{0.23\textwidth}|}
\hline
\textbf{Endpoint} & \textbf{Кто вызывает} & \textbf{Назначение} \\
\hline
\pathcode{GET /identity/internal/v1/auth/jwks} & API Gateway и сервисы & Получение публичных ключей для проверки JWT \\
\hline
\pathcode{GET /identity/internal/v1/users/\{userId\}/summary} & Review Service, Booking Service & Получение краткого профиля пользователя \\
\hline
\pathcode{GET /catalog/internal/v1/restaurants/\{restaurantId\}/summary} & Booking Service, Review Service & Получение краткой информации о ресторане \\
\hline
\pathcode{GET /catalog/internal/v1/restaurants/\{restaurantId\}/booking-context} & Booking Service & Проверка ресторана, столика, вместимости и рабочих часов перед созданием бронирования \\
\hline
\pathcode{GET /catalog/internal/v1/restaurants/\{restaurantId\}/availability-context} & Booking Service & Получение активных столиков и рабочих часов для расчета доступности \\
\hline
\pathcode{POST /booking/internal/v1/bookings/availability/check} & Restaurant Catalog Service или API Gateway & Проверка занятости набора столиков на интервал времени \\
\hline
\pathcode{GET /booking/internal/v1/bookings/\{bookingId\}/review-context} & Review Service & Проверка права пользователя оставить отзыв \\
\hline
\pathcode{GET /reviews/internal/v1/restaurants/\{restaurantId\}/rating-stats} & Restaurant Catalog Service & Синхронизация рейтингов при пересчете или восстановлении read-model \\
\hline
\end{longtable}

Для внутренних запросов используется отдельная сервисная авторизация, например заголовок \code{X-Service-Token}. Пользовательский JWT не должен использоваться как единственный механизм доверия между сервисами.

Асинхронное событие \code{ReviewCreated} передается из Review Service в Catalog Service после создания отзыва. Минимальный payload события:

\begin{lstlisting}[style=jsonstyle]
{
  "eventId": "9d15a4f0-4c47-4d68-b79a-6bb0d1e2a111",
  "eventType": "ReviewCreated",
  "occurredAt": "2026-05-12T12:00:00Z",
  "reviewId": 101,
  "restaurantId": 1,
  "userId": 3,
  "rating": 5
}
\end{lstlisting}

Catalog Service использует событие для пересчета \code{restaurant\_rating\_stats}. Повторная доставка события должна быть безопасной: для этого можно хранить обработанные \code{eventId} или пересчитывать агрегат по данным Review Service через internal endpoint.

\subsection{Основные сценарии взаимодействия}

\subsubsection{Создание бронирования}

\begin{enumerate}[leftmargin=1.7cm]
    \item Клиент отправляет \pathcode{POST /api/v1/bookings} в API Gateway.
    \item Gateway проверяет JWT и передает запрос в Booking Service вместе с \code{userId}.
    \item Booking Service вызывает \pathcode{GET /catalog/internal/v1/restaurants/\{restaurantId\}/booking-context}.
    \item Catalog Service проверяет, что ресторан и столик активны, столик принадлежит ресторану, вместимость подходит, а интервал попадает в рабочие часы.
    \item Booking Service проверяет в \code{booking\_db}, что у выбранного столика нет пересекающегося активного бронирования.
    \item Booking Service создает запись в \code{bookings} и сохраняет snapshot ресторана и столика.
    \item Клиент получает созданное бронирование.
\end{enumerate}

Возможные ошибки: \code{401} при отсутствии JWT, \code{404} если ресторан или столик не найдены, \code{422} если столик не подходит по вместимости или времени работы, \code{409} если столик уже забронирован на выбранное время.

\subsubsection{Получение доступности ресторана}

\begin{enumerate}[leftmargin=1.7cm]
    \item Клиент отправляет \pathcode{GET /api/v1/restaurants/\{restaurantId\}/availability}.
    \item Gateway маршрутизирует запрос в Booking Service.
    \item Booking Service получает из Catalog Service рабочие часы и активные столики.
    \item Booking Service проверяет пересечения с бронированиями в своей БД.
    \item Клиент получает список свободных слотов.
\end{enumerate}

Такой подход сохраняет единственного владельца данных о бронированиях и убирает прямой SQL join между \code{restaurant\_tables} и \code{bookings}.

\subsubsection{Создание отзыва}

\begin{enumerate}[leftmargin=1.7cm]
    \item Клиент отправляет \pathcode{POST /api/v1/restaurants/\{restaurantId\}/reviews}.
    \item Gateway проверяет JWT и передает запрос в Review Service.
    \item Review Service вызывает \pathcode{GET /booking/internal/v1/bookings/\{bookingId\}/review-context}.
    \item Booking Service проверяет, что бронирование принадлежит пользователю, относится к выбранному ресторану, завершено и уже прошло по времени.
    \item Review Service проверяет уникальность \code{booking\_id} в своей БД.
    \item Review Service создает отзыв и записывает событие \code{ReviewCreated} в \code{review\_outbox}.
    \item Catalog Service получает событие и обновляет \code{restaurant\_rating\_stats}.
\end{enumerate}

Возможные ошибки: \code{404} если бронирование не найдено или не принадлежит пользователю, \code{400} если оценка вне диапазона, \code{409} если отзыв по бронированию уже существует, \code{422} если бронирование еще не завершено.

\subsubsection{Поиск ресторанов с сортировкой по рейтингу}

\begin{enumerate}[leftmargin=1.7cm]
    \item Клиент отправляет \pathcode{GET /api/v1/restaurants} с фильтрами, сортировкой и пагинацией.
    \item Gateway маршрутизирует запрос в Restaurant Catalog Service.
    \item Catalog Service фильтрует рестораны по своим таблицам и использует локальную таблицу \code{restaurant\_rating\_stats}.
    \item Если read-model нужно восстановить, Catalog Service может запросить \pathcode{GET /reviews/internal/v1/restaurants/\{restaurantId\}/rating-stats}.
\end{enumerate}

Поиск не должен делать синхронный запрос в Review Service для каждого ресторана из страницы, иначе пагинация и сортировка станут нестабильными и медленными.

\subsection{Общий формат ошибок}

Для публичных и внутренних endpoint-ов используется единый формат ошибки:

\begin{lstlisting}[style=jsonstyle]
{
  "timestamp": "2026-05-12T12:00:00Z",
  "status": 409,
  "error": "Conflict",
  "message": "This table is already booked for the selected time range",
  "path": "/api/v1/bookings",
  "details": []
}
\end{lstlisting}

Основные коды:

\begin{longtable}{|L{0.25\textwidth}|L{0.65\textwidth}|}
\hline
\textbf{Код} & \textbf{Когда возвращается} \\
\hline
\code{400 Bad Request} & Ошибка формата запроса или валидации DTO \\
\hline
\code{401 Unauthorized} & Нет JWT или service token \\
\hline
\code{403 Forbidden} & Пользователь или сервис не имеет прав \\
\hline
\code{404 Not Found} & Сущность не найдена или недоступна вызывающей стороне \\
\hline
\code{409 Conflict} & Конфликт текущего состояния: занятый столик, повторный отзыв \\
\hline
\code{422 Unprocessable Entity} & Запрос синтаксически корректен, но нарушает бизнес-правило \\
\hline
\code{503 Service Unavailable} & Зависимый сервис временно недоступен \\
\hline
\end{longtable}

\subsection{План миграции монолита}

\begin{enumerate}[leftmargin=1.7cm]
    \item Зафиксировать текущий внешний REST API из ДЗ2 и ЛР1 как контракт, который должен сохранить API Gateway.
    \item Вынести Identity Service: перенести таблицу \code{users}, регистрацию, логин, профиль пользователя и выпуск JWT.
    \item Перевести JWT на схему, удобную для микросервисов: Identity Service выпускает токены, Gateway и сервисы проверяют подпись по публичному ключу или по общему internal contract.
    \item Вынести Restaurant Catalog Service: перенести таблицы ресторанов, кухонь, столиков, меню, фото и рабочих часов.
    \item Добавить в Catalog Service internal endpoint-ы для получения restaurant/table context, чтобы Booking Service не обращался к чужой БД.
    \item Вынести Booking Service: перенести \code{bookings}, убрать внешние ключи на чужие таблицы, добавить snapshot-поля для отображения бронирований.
    \item Перенести endpoint availability в Booking Service, потому что расчет зависит от бронирований.
    \item Вынести Review Service: перенести \code{reviews}, убрать внешние ключи на \code{bookings}, \code{users}, \code{restaurants}, оставить scalar id и unique constraint на \code{booking\_id}.
    \item Добавить проверку права на отзыв через internal endpoint Booking Service.
    \item Добавить outbox в Review Service и обработчик события \code{ReviewCreated} в Catalog Service для обновления \code{restaurant\_rating\_stats}.
    \item Подготовить отдельные Liquibase changelog-и для каждой БД.
    \item Перенести данные из монолитной схемы \code{storage} в новые базы: \code{identity\_db}, \code{catalog\_db}, \code{booking\_db}, \code{review\_db}.
    \item Выполнить regression smoke-test по Postman-сценарию из ДЗ3: login, users/me, cuisines, restaurant search, details, menu, availability, create booking, my bookings, cancel booking, reviews.
    \item После стабилизации выключить прямые обращения монолита к старой общей БД.
\end{enumerate}

\subsection{Требования к надежности}

Распределенные транзакции между сервисами не используются. Если операция затрагивает несколько сервисов, владелец основной бизнес-операции сохраняет свое состояние локально, а дополнительные изменения выполняются через события или повторяемые internal API.

Для событий используется outbox-подход: Review Service сначала сохраняет отзыв и событие в одной локальной транзакции, затем отдельный publisher отправляет событие в брокер. Catalog Service должен обрабатывать событие идемпотентно, чтобы повторная доставка не ломала статистику рейтинга.

Для диагностики каждый входящий запрос получает correlation id. Gateway передает его во все downstream-сервисы, а сервисы пишут его в логи.

\section{Вывод}

В результате был подготовлен технический дизайн микросервисной архитектуры для приложения бронирования столиков в ресторанах. Монолит разделен на четыре предметных сервиса и API Gateway. Для каждого сервиса определены зона ответственности и собственная база данных. Прямые связи между таблицами заменены внутренними REST-контрактами и асинхронным событием обновления рейтинга. Также подготовлена OpenAPI-спецификация межсервисных endpoint-ов и описан пошаговый план миграции от текущего монолита к архитектуре \code{database-per-service}.

\end{document}
